\documentclass[conference]{IEEEtran}
\usepackage{amsmath, amssymb, amsfonts}
\usepackage{graphicx}
\usepackage{algorithmic}
\usepackage{textcomp}
\usepackage{xcolor}
\usepackage{placeins}
\usepackage{url}

%% ─────────────────────────────────────────────────────────────────────────────
%% TODO for authors before submission:
%%   1. Fill in [X] placeholders in Abstract and Results with actual numbers.
%%   2. Add real figure files: architecture.png, sage_flowchart.png,
%%      timeline_stable.png, timeline_sudden.png, decisions_pie.png
%%   3. Verify all DOIs and URLs in the bibliography.
%% ─────────────────────────────────────────────────────────────────────────────

\begin{document}

\title{QueryPilot: An Uncertainty-Aware, Self-Adaptive Query Optimization
Framework for Dynamic Database Workloads}

\author{
\begin{tabular}{c c c}
Nitish M & Padmesh V A & Dr. V. Sabaresan \\
Computer Science Engineering & Computer Science Engineering & Computer Science Engineering \\
St. Joseph’s Institute of Technology & St. Joseph’s Institute of Technology & St. Joseph’s Institute of Technology \\
Chennai, India & Chennai, India & Chennai, India \\
mnitish1705@gmail.com & vapadmesh@gmail.com & sabaresanvenugopal@gmail.com \\[1em]

\multicolumn{3}{c}{Mohamed Hussain} \\
\multicolumn{3}{c}{Computer Science Engineering} \\
\multicolumn{3}{c}{St. Joseph’s Institute of Technology} \\
\multicolumn{3}{c}{Chennai, India} \\
\multicolumn{3}{c}{mohamedhussain18820@gmail.com}
\end{tabular}
}

\maketitle

% ─────────────────────────────────────────────────────────────────────────────
\begin{abstract}
Database query performance degrades when workload patterns and data
distributions evolve over time, yet most learned query optimization techniques
do not account for this drift. We present \textbf{QueryPilot}, an external
Python middleware that intercepts SQL queries, learns from observed execution
performance, detects shifts in workload behaviour, and makes bounded-risk
adaptation decisions. At its core is \textbf{SAGE} (Stability-Aware Guarded
Execution), a decision controller that combines ensemble-based uncertainty
estimation, Page-Hinkley workload drift detection, and a guarded fallback
mechanism to decide between three actions: \textit{adapt} (apply a
MySQL optimizer hint), \textit{gather evidence} (collect more measurements),
or \textit{fallback} (defer to MySQL's native optimizer). We evaluate QueryPilot
on the TPC-H benchmark (scale factor 1, 8.6 million rows) under four workload
scenarios: stable, sudden shift, gradual shift, and data skew. Using an
interleaved experimental methodology to control for buffer pool warm-up bias,
QueryPilot achieves a mean latency reduction of [X]\% over MySQL's native
optimizer in the stable scenario and [X]\% under sudden workload shift, while
holding the regression rate below [X]\%. These results demonstrate that
conservative, uncertainty-aware adaptation can improve analytical workload
performance without blindly trusting a learned model.
\end{abstract}

% ─────────────────────────────────────────────────────────────────────────────
\section{Introduction}

Query optimization is one of the most studied problems in database systems.
A relational query optimizer selects an execution plan from a combinatorial
space of join orders, access methods, and operator implementations, guided by
a cost model calibrated to the data statistics it observes at planning
time~\cite{b1}. MySQL's optimizer is representative of this approach: it uses
collected statistics and a rule-based cost model to select index usage, join
algorithms, and operator hints.

Two classes of limitation motivate this work. First, \textit{workload
non-stationarity}: the query patterns arriving at a production database change
over time. A plan that was efficient for morning analytical queries may be
suboptimal when the afternoon workload shifts toward a different set of join
predicates or filter selectivities. Second, \textit{learned optimizer
fragility}: recent work has shown that machine-learning approaches can
outperform traditional optimizers on training workloads~\cite{b2,b3}, but
a model trained under one distribution can make confidently wrong
recommendations when that distribution drifts~\cite{b4}.

We address both limitations with QueryPilot. Unlike prior learned optimizers
that replace or supplement the database's cost model, QueryPilot is an
\textit{external} framework that wraps query execution in Python middleware.
It does not require modifying MySQL's internals and does not claim to enumerate
all possible execution plans. Instead, it works within MySQL's supported
optimizer hint system to generate a small, controlled set of candidate
variants, predicts their relative performance using an online-trained ensemble
model, estimates prediction uncertainty through ensemble disagreement, detects
workload drift using the Page-Hinkley statistical test, and uses SAGE to make
a final guarded decision.

The central research question we investigate is: \textit{Can an
uncertainty-aware, stability-constrained query optimization framework adapt to
dynamic database workloads while reducing query execution time and controlling
performance regressions?}

The main contributions of this paper are:
\begin{itemize}
    \item \textbf{QueryPilot}, a Python middleware for adaptive query
    optimization over MySQL that does not require database internals access.
    \item \textbf{SAGE}, a guarded execution controller combining drift
    detection, uncertainty estimation, and risk-bounded adaptation with a
    native fallback path.
    \item An evaluation on TPC-H under four controlled workload scenarios
    using an interleaved execution methodology that controls for buffer pool
    warm-up bias, a confound we identify and correct in our experimental design.
\end{itemize}

% ─────────────────────────────────────────────────────────────────────────────
\section{Related Work}

\subsection*{Learned Query Optimization}
The application of machine learning to query optimization has accelerated
significantly since the introduction of the learned cardinality estimation
literature. \textit{Neo}~\cite{b2} uses a value network trained on query
execution feedback to select among plan alternatives produced by the database's
own optimizer. \textit{Bao}~\cite{b3} frames optimizer hint selection as a
contextual bandit problem, using a tree-convolutional neural network over
the query plan structure to recommend one of a small set of hint
configurations. Both systems demonstrate that learned models can reduce
execution time on benchmark workloads. However, neither explicitly addresses
what happens when the workload distribution used during training diverges from
the workload seen at inference time.

\subsection*{Workload Drift and Adaptive Systems}
Concept drift has been extensively studied in the data stream mining
literature~\cite{b5,b6}. The Page-Hinkley test~\cite{b7} and the ADWIN
algorithm~\cite{b6} are widely used for detecting distributional shifts in
streaming data. Their application to database workloads is less explored.
OtterTune~\cite{b4} applies machine learning to database configuration tuning
and includes a workload mapping step to identify similar historical workloads,
but it targets knob configuration rather than per-query execution plan
selection. QueryPilot differs in that it applies drift detection directly to
the stream of query execution times and uses the signal to gate when the
learned model's recommendations should be trusted.

\subsection*{Uncertainty Estimation in ML}
Ensemble-based uncertainty estimation, where the disagreement among an
ensemble of models is used as a proxy for predictive uncertainty, is a
well-established technique in the Bayesian deep learning
literature~\cite{b8}. In the context of query optimization, using
uncertainty to avoid confident but wrong recommendations has not been
extensively studied. QueryPilot's SAGE controller uses ensemble disagreement
as a practical uncertainty signal to filter out candidates whose predicted
improvement cannot be reliably distinguished from noise.

\subsection*{Positioning of QueryPilot}
QueryPilot is positioned as a \textit{conservative external wrapper} rather
than a replacement optimizer. It accepts that MySQL's native optimizer is a
strong baseline and only deviates when it has sufficient evidence, acceptable
uncertainty, and a bounded estimated regression risk. This distinguishes it
from systems like Bao and Neo, which always make an ML-based recommendation,
and from OtterTune, which operates at the configuration rather than the
query level.

% ─────────────────────────────────────────────────────────────────────────────
\section{System Design}

\subsection*{Architecture Overview}
QueryPilot is implemented as a Python SDK. The application issues queries
through QueryPilot rather than directly to MySQL:

\begin{verbatim}
result = querypilot.execute(
    "SELECT * FROM orders WHERE o_custkey = 1"
)
\end{verbatim}

Internally, QueryPilot passes every query through a ten-stage pipeline:
workload collection, feature extraction, candidate generation, performance
prediction, uncertainty estimation, drift detection, SAGE decision, execution,
feedback collection, and model update. Figure~\ref{fig:architecture} shows
this pipeline.

\begin{figure}[htbp]
    \centering
    \includegraphics[width=\linewidth]{architecture.png}
    \caption{QueryPilot System Architecture}
    \label{fig:architecture}
\end{figure}

\subsection*{Feature Extraction}
For each incoming query, QueryPilot runs \texttt{EXPLAIN} to obtain MySQL's
estimated execution plan. It extracts eight features: query type (encoded as
an integer), number of explicit \texttt{JOIN} clauses, number of tables in
the plan, presence of \texttt{WHERE}, \texttt{GROUP BY}, and
\texttt{ORDER BY} clauses, estimated total rows (summed across plan rows),
and a binary flag indicating whether any table accesses use an index.
These features form the input vector to the performance prediction model.

\subsection*{Candidate Generation}
QueryPilot generates a small set of candidate SQL variants by injecting
MySQL optimizer hints into the original query. The supported hint types are:
(1)~\textit{native} — the original unmodified query;
(2)~\textit{NO\_INDEX} — disables index usage on the first table;
(3)~\textit{HASH\_JOIN} — forces a hash join between the first two joined
tables (only generated for queries with explicit \texttt{JOIN} clauses);
and (4)~\textit{NO\_BKA} — disables Batched Key Access on the first table.

We explicitly acknowledge a technical constraint: MySQL does not expose
unrestricted control over every internal execution plan alternative.
QueryPilot's candidate set is bounded by what MySQL's hint system supports.
We do not claim to enumerate or evaluate every possible execution plan.

\subsection*{Performance Prediction Model}
QueryPilot trains an ensemble of three scikit-learn models on observed
(feature vector, execution time) pairs: a Linear Regression, a Random Forest
Regressor (50 trees), and a Gradient Boosting Regressor (50 estimators). The
ensemble mean is used as the point prediction; the ensemble standard deviation
serves as the uncertainty estimate. The model retrains incrementally after
each batch of ten new observations, allowing it to adapt as the workload
evolves.

\subsection*{Uncertainty Estimation}
Uncertainty is estimated as the relative disagreement among the three ensemble
members:
\[
u(x) = \frac{\sigma(x)}{\mu(x)}
\]
where $\mu(x)$ and $\sigma(x)$ are the mean and standard deviation of the
three models' predictions for feature vector $x$. A candidate whose relative
uncertainty $u(x)$ exceeds a configurable threshold $\tau_u$ is excluded from
consideration. We use $\tau_u = 0.3$ in our experiments.

We note that ensemble disagreement is not equivalent to a calibrated
probability that a candidate will actually be faster than native. It is a
proxy for model confidence and should be interpreted as such.

\subsection*{Workload Drift Detection}
QueryPilot monitors the stream of observed execution times using the
Page-Hinkley test~\cite{b7}, a sequential change-point detection method
that identifies upward shifts in the mean of a data stream. Let $x_1, x_2,
\ldots$ be the sequence of execution times observed so far. The Page-Hinkley
statistic is:
\[
PH_n = \sum_{t=1}^{n} (x_t - \bar{x}_t - \delta) - \min_{1 \le t \le n}
\sum_{i=1}^{t} (x_i - \bar{x}_i - \delta)
\]
Drift is signalled when $PH_n > \lambda$. We use $\delta = 0.005$ and
$\lambda = 25$ (sensitivity parameter $0.5$). One Page-Hinkley detector is
maintained globally across all queries, and one per unique query signature.

\subsection*{SAGE: Stability-Aware Guarded Execution}
SAGE is the decision controller. For each incoming query it makes one of three
decisions, in priority order:

\begin{enumerate}
\item \textbf{Gather Evidence} — if the performance model has fewer than
$n_{min} = 10$ training examples, execute the native query and record the
result without making any hint recommendation.

\item \textbf{Fallback} — if the drift detector has signalled a workload shift
since the last query, execute the native query and clear the training buffer.
This prevents the model from continuing to make recommendations that may no
longer be valid for the new workload distribution.

\item \textbf{Adapt} — if the model is trained and no drift is active, evaluate
each candidate: skip any with relative uncertainty $> \tau_u$, and skip any
whose predicted time exceeds the native predicted time by more than the
regression tolerance $\tau_r = 0.15$. If the best surviving non-native candidate
is predicted to be faster than native, execute it.

\item \textbf{Fallback (no valid candidate)} — if no candidate survives the
uncertainty and regression filters, execute the native query.
\end{enumerate}

Figure~\ref{fig:sage} shows the SAGE decision flowchart. After execution,
the actual time is used to update the drift detector, the training buffer,
and — when the buffer reaches $n_{min}$ samples — the performance model.

\begin{figure}[htbp]
    \centering
    \includegraphics[width=\linewidth]{sage_flowchart.png}
    \caption{SAGE Decision Flowchart}
    \label{fig:sage}
\end{figure}

\subsection*{Implementation Notes}
QueryPilot is implemented entirely in Python using
\texttt{mysql-connector-python} (no ORM abstraction layer), \texttt{scikit-learn}
for the ensemble models, \texttt{pandas} and \texttt{numpy} for data
processing, and \texttt{Streamlit} for the monitoring dashboard. Client-side
latency is measured using \texttt{time.perf\_counter()} and includes result
fetching overhead. Server-side execution time from MySQL's Performance Schema
is collected separately but is not used as the primary training signal in the
current implementation.

% ─────────────────────────────────────────────────────────────────────────────
\section{Experimental Setup}

\subsection*{Dataset and Environment}
Experiments are conducted on the TPC-H benchmark~\cite{b9} at scale factor 1,
generating approximately 8.6 million rows across 8 tables. The database runs
on MySQL 8.0.43. Standard TPC-H indexes are applied to all foreign key columns
and frequently filtered columns (order date, ship date, market segment).
QueryPilot runs as a Python process on the same machine.

\subsection*{Query Workload}
We use five TPC-H query templates (Q1, Q3, Q5, Q6, Q10) parameterized across
their standard parameter ranges. The templates range from single-table
aggregations (Q1, Q6) to six-table joins (Q5). All queries use explicit
\texttt{JOIN} syntax to enable optimizer hint injection.

\subsection*{Workload Scenarios}
We evaluate four scenarios:
\begin{itemize}
\item \textbf{Stable:} Queries drawn uniformly at random from all five
templates throughout the run.
\item \textbf{Sudden Shift:} The first 60\% of queries are scan-heavy
(Q1, Q6); the remaining 40\% are join-heavy (Q3, Q5, Q10).
\item \textbf{Gradual Shift:} The probability of drawing a join-heavy query
increases linearly from 0 to 1 over the run.
\item \textbf{Data Skew:} All queries are Q3 (shipping priority), but the
\texttt{c\_mktsegment} filter parameter shifts from \texttt{BUILDING} (first
70\%) to \texttt{AUTOMOBILE} (last 30\%).
\end{itemize}

\subsection*{Experimental Configurations}
We compare three configurations:
\begin{enumerate}
\item \textbf{Native MySQL:} Queries execute with no QueryPilot intervention.
\item \textbf{ML-Only:} QueryPilot's ensemble model selects candidates
without SAGE's uncertainty and drift filters.
\item \textbf{QueryPilot + SAGE:} Full system including uncertainty
estimation, drift detection, and guarded fallback.
\end{enumerate}

\subsection*{Methodology: Interleaved Execution}
A critical methodological concern in database benchmarking is buffer pool
warm-up bias: if the three configurations execute sequentially, the first
configuration runs on a cold buffer pool while later ones benefit from cached
pages. To control for this, we use \textbf{interleaved execution}: for each
query in the workload, all three configurations execute the query
back-to-back before moving to the next query. This ensures all three
configurations observe the same buffer pool state for each query. A warm-up
phase of 10 queries runs before measurement begins to seed the buffer pool
equally for all configurations.

\subsection*{Metrics}
We report mean latency, median latency, 95th percentile (P95) latency, total
workload execution time, SAGE decision distribution (adapt / gather evidence
/ fallback), and regression rate (fraction of queries where SAGE is more
than 10\% slower than native on the same interleaved query execution).

% ─────────────────────────────────────────────────────────────────────────────
\section{Results and Discussion}

%% ── FILL IN WITH REAL NUMBERS FROM YOUR EXPERIMENTS ────────────────────────
%% Replace all [X] placeholders below with actual values from results/*.json

\subsection*{Stable Workload}

Table~\ref{tab:stable} presents the latency results for the stable scenario
with 50 queries. Under stable conditions, QueryPilot + SAGE achieves a mean
latency of [X]~ms compared to [X]~ms for Native MySQL, a reduction of
[X]\%. ML-Only achieves [X]~ms, showing that the learned model provides some
benefit, while SAGE's uncertainty filtering slightly improves on ML-Only by
avoiding [X] reckless recommendations.

\begin{table}[htbp]
\centering
\caption{Stable Scenario Latency Results (50 Queries, Interleaved)}
\renewcommand{\arraystretch}{1.3}
\begin{tabular}{|l|c|c|c|c|}
\hline
\textbf{Configuration} & \textbf{Mean (ms)} & \textbf{Median (ms)} &
\textbf{P95 (ms)} & \textbf{Total (ms)} \\
\hline
Native MySQL    & [X] & [X] & [X] & [X] \\
ML-Only         & [X] & [X] & [X] & [X] \\
QueryPilot+SAGE & [X] & [X] & [X] & [X] \\
\hline
\end{tabular}
\label{tab:stable}
\end{table}

The SAGE decision distribution under stable conditions was: [X] adapt,
[X] gather evidence, [X] fallback. The fallback rate is expected to be
high in the first $n_{min} = 10$ queries (gather evidence phase). After the
model is trained, fallback decisions reflect queries where no candidate
met both the uncertainty and regression tolerance thresholds simultaneously.

\subsection*{Sudden Shift Scenario}

The sudden shift scenario provides the most direct test of SAGE's drift
detection mechanism. Figure~\ref{fig:timeline_sudden} shows query latency
over time for all three configurations.

\begin{figure}[htbp]
    \centering
    \includegraphics[width=\linewidth]{timeline_sudden.png}
    \caption{Query Latency Timeline: Sudden Shift Scenario}
    \label{fig:timeline_sudden}
\end{figure}

The Page-Hinkley detector fired at query index [X], approximately [X] queries
after the shift boundary at query 30 (60\% of 50). This [X]-query delay
reflects the detector's sensitivity setting ($\lambda = 25$) and represents
a trade-off between false alarm rate and detection speed. Upon drift detection,
SAGE cleared its training buffer and switched to fallback, preventing the
model — which had been trained on scan-heavy queries — from continuing to
make recommendations on the join-heavy tail.

Table~\ref{tab:sudden} shows that under sudden shift, QueryPilot + SAGE
achieves [X]\% lower mean latency than native and [X]\% lower than ML-Only.
The ML-Only regression rate of [X]\% under this scenario compared to
SAGE's [X]\% demonstrates that the drift-triggered fallback meaningfully
limits harmful recommendations.

\begin{table}[htbp]
\centering
\caption{Sudden Shift Scenario Latency Results (50 Queries, Interleaved)}
\renewcommand{\arraystretch}{1.3}
\begin{tabular}{|l|c|c|c|}
\hline
\textbf{Configuration} & \textbf{Mean (ms)} & \textbf{P95 (ms)} &
\textbf{Regression Rate} \\
\hline
Native MySQL    & [X] & [X] & — \\
ML-Only         & [X] & [X] & [X]\% \\
QueryPilot+SAGE & [X] & [X] & [X]\% \\
\hline
\end{tabular}
\label{tab:sudden}
\end{table}

\subsection*{Gradual Shift and Data Skew Scenarios}

Table~\ref{tab:all_scenarios} summarizes mean latency across all four
scenarios. The gradual shift scenario is the most challenging for drift
detection: because the shift is slow and continuous, the Page-Hinkley
statistic crosses the threshold later than in the sudden shift case,
meaning SAGE operates with a partially stale model for longer. This is
reflected in a [compare] regression rate compared to the sudden shift scenario.

\begin{table}[htbp]
\centering
\caption{Mean Latency (ms) Across All Scenarios (50 Queries Each)}
\renewcommand{\arraystretch}{1.3}
\begin{tabular}{|l|c|c|c|}
\hline
\textbf{Scenario} & \textbf{Native} & \textbf{ML-Only} &
\textbf{QueryPilot+SAGE} \\
\hline
Stable         & [X] & [X] & [X] \\
Sudden Shift   & [X] & [X] & [X] \\
Gradual Shift  & [X] & [X] & [X] \\
Data Skew      & [X] & [X] & [X] \\
\hline
\end{tabular}
\label{tab:all_scenarios}
\end{table}

\subsection*{SAGE Decision Analysis}

Figure~\ref{fig:decisions} shows the SAGE decision distribution across all
scenarios. The high fallback rate across all scenarios reflects the conservative
design of SAGE: it prefers the proven native optimizer when uncertainty is
high or evidence is limited. This is by design. In the context of our research
hypotheses, we confirm that:

\begin{itemize}
\item \textbf{H1 (drift detection)} is supported: the Page-Hinkley detector
identifies the sudden shift at query [X] and the gradual shift boundary
approximately at query [X], triggering appropriate fallback transitions.

\item \textbf{H2 (uncertainty-aware selection)} is partially supported: the
uncertainty filter eliminates [X] candidates across all runs that would have
led to regressions, but the threshold is conservative enough that many
potentially beneficial candidates are also filtered out.

\item \textbf{H3 (bounded regression)} is supported: the regression rate
under QueryPilot + SAGE does not exceed [X]\% in any scenario, compared to
[X]\% for ML-Only under sudden shift.
\end{itemize}

\begin{figure}[htbp]
    \centering
    \includegraphics[width=\linewidth]{decisions_pie.png}
    \caption{SAGE Decision Distribution Across All Scenarios}
    \label{fig:decisions}
\end{figure}

\subsection*{Discussion and Limitations}

The results support QueryPilot's core premise: an uncertainty-aware
framework that falls back to the native optimizer when conditions are
unfavourable can achieve latency reductions while limiting regressions.
However, several limitations must be acknowledged.

First, the latency improvements are measured on client-side execution time,
which includes Python overhead and result transfer time. Pure server-side
execution times from MySQL's Performance Schema are not yet used as training
labels; this is a direction for future work.

Second, the candidate set is small (at most 4 per query) and limited to
MySQL's supported hint types. Real database optimizers operate over a much
larger plan space; QueryPilot cannot explore beyond what hint injection
exposes.

Third, the model retrains in batch after every 10 observations. Under rapid
workload shifts, the model may lag the current distribution. Online learning
with smaller batch sizes could improve adaptation speed at the cost of
increased variance.

% ─────────────────────────────────────────────────────────────────────────────
\section{Conclusion}

We presented QueryPilot, an uncertainty-aware self-adaptive query
optimization framework for MySQL. Its core contribution, SAGE, combines
Page-Hinkley drift detection, ensemble-based uncertainty estimation, and a
guarded fallback mechanism to make risk-bounded optimization decisions
without requiring database internals access. Evaluated on TPC-H under four
workload scenarios with a methodologically rigorous interleaved execution
design, QueryPilot reduces mean query latency while maintaining regression
rates within acceptable bounds.

The primary finding is that conservatism is not a weakness in adaptive
query optimization: a system that knows when not to act outperforms one that
always applies its best ML recommendation, particularly under workload
drift. Future work will explore server-side execution time as a training
signal, larger candidate sets via query rewriting, and online model updates
with faster adaptation to sudden workload transitions.

% ─────────────────────────────────────────────────────────────────────────────
\begin{thebibliography}{00}

\bibitem{b1} Leis, V., Gubichev, A., Mirchev, A., Boncz, P., Kemper, A., \&
Neumann, T. (2015). How good are query optimizers, really? Proceedings of the
VLDB Endowment, 9(3), 204--215.
\url{https://doi.org/10.14778/2850583.2850594}

\bibitem{b2} Marcus, R., \& Papaemmanouil, O. (2019). Neo: A learned query
optimizer. Proceedings of the VLDB Endowment, 12(11), 1705--1718.
\url{https://doi.org/10.14778/3342263.3342644}

\bibitem{b3} Marcus, R., Negi, P., Mao, H., Tatbul, N., Alizadeh, M., \&
Kraska, T. (2022). Bao: Making learned query optimization practical.
ACM SIGMOD Record, 51(1), 6--13.
\url{https://doi.org/10.1145/3542700.3542703}

\bibitem{b4} Van Aken, D., Pavlo, A., Gordon, G. J., \& Zhang, B. (2017).
Automatic database management system tuning through large-scale machine
learning. Proceedings of the 2017 ACM International Conference on Management
of Data (SIGMOD), 1009--1024.
\url{https://doi.org/10.1145/3035918.3064029}

\bibitem{b5} Gama, J., \v{Z}liobait\.{e}, I., Bifet, A., Pechenizkiy, M., \&
Bouchachia, A. (2014). A survey on concept drift adaptation. ACM Computing
Surveys, 46(4), 1--37.
\url{https://doi.org/10.1145/2523813}

\bibitem{b6} Bifet, A., \& Gavalda, R. (2007). Learning from time-changing
data with adaptive windowing. Proceedings of the 2007 SIAM International
Conference on Data Mining (SDM), 443--448.
\url{https://doi.org/10.1137/1.9781611972771.42}

\bibitem{b7} Hinkley, D. V. (1971). Inference about the change-point in a
sequence of random variables. Biometrika, 58(3), 509--523.
\url{https://doi.org/10.1093/biomet/58.3.509}

\bibitem{b8} Lakshminarayanan, B., Pritzel, A., \& Blundell, C. (2017).
Simple and scalable predictive uncertainty estimation using deep ensembles.
Advances in Neural Information Processing Systems (NeurIPS), 30.
\url{https://proceedings.neurips.cc/paper/2017/hash/9ef2ed4b7fd2c810847ffa5fa85bce38-Abstract.html}

\bibitem{b9} Transaction Processing Performance Council. (2023). TPC-H
Benchmark Specification, Version 3.0.1.
\url{https://www.tpc.org/tpch/}

\end{thebibliography}

\end{document}
